% Part: methods
% Chapter: proofs
% Section: proof-by-contradiction

\documentclass[../../../include/open-logic-section]{subfiles}

\begin{document}

\olfileid[id]{mth}{prf}{con}

\olsection{Pembuktian dengan Kontradiksi}

Pada mulanya, pembuktian dengan kontradiksi merupakan pola inferensi
yang digunakan untuk membuktikan klaim negatif. Andaikan Anda ingin
menunjukkan bahwa suatu klaim~$p$ \emph{salah}, yaitu Anda ingin
menunjukkan~$\lnot p$. Strategi yang paling menjanjikan ialah (a)
mengandaikan bahwa $p$~benar, dan (b) menunjukkan bahwa asumsi ini
menuntun pada sesuatu yang Anda ketahui salah. ``Sesuatu yang diketahui
salah'' dapat berupa hasil yang bertentangan---berkontradiksi---dengan
$p$ itu sendiri, atau dengan hipotesis lain dari keseluruhan klaim
yang sedang Anda telaah. Misalnya, bukti atas ``jika $q$, maka $\lnot
p$'' melibatkan pengasumsian bahwa $q$~benar dan pembuktian~$\lnot p$
darinya. Jika Anda membuktikan $\lnot p$ dengan kontradiksi, itu
berarti mengasumsikan $p$ di samping~$q$. Jika Anda dapat membuktikan
$\lnot q$ dari $p$, Anda telah menunjukkan bahwa asumsi~$p$ menuntun
pada sesuatu yang bertentangan dengan asumsi Anda yang lain, yaitu
~$q$, sebab $q$~dan $\lnot q$ tidak dapat sama-sama benar. Tentu saja,
dalam membuktikan kontradiksi Anda harus menggunakan pola inferensi
lain serta menguraikan definisi. Mari kita telaah sebuah contoh.

\begin{prop}
  Jika $A \subseteq B$ dan $B = \emptyset$, maka $A$ tidak memiliki
  !!{element}.
\end{prop}

\begin{proof}
  Andaikan $A \subseteq B$ dan $B = \emptyset$. Kita ingin menunjukkan
  bahwa $A$ tidak memiliki !!{element}.
  \begin{quote}
    Karena ini merupakan klaim kondisional, kita mengasumsikan
    anteseden dan ingin membuktikan konsekuennya. Konsekuennya ialah:
    $A$ tidak memiliki !!{element}. Kita dapat membuatnya sedikit
    lebih eksplisit: tidak benar bahwa ada~$x \in A$.
  \end{quote}
  $A$ tidak memiliki !!{element} jika dan hanya jika tidak benar bahwa
  ada~$x$ sedemikian sehingga $x \in A$.
  \begin{quote}
    Jadi, kita telah menentukan bahwa yang sebenarnya ingin kita
    buktikan adalah klaim negatif~$\lnot p$, yaitu: tidak benar bahwa
    ada $x \in A$. Untuk menggunakan pembuktian dengan kontradiksi,
    kita harus mengasumsikan klaim positif yang bersesuaian, yaitu~$p$:
    ada $x \in A$, lalu membuktikan suatu kontradiksi darinya. Kita
    menandai bahwa kita sedang melakukan pembuktian dengan kontradiksi
    dengan menulis ``untuk memperoleh kontradiksi, asumsikan'' atau
    bahkan sekadar ``andaikan tidak'', lalu menyatakan asumsi~$p$.
  \end{quote}
  Andaikan tidak: ada $x \in A$.
  \begin{quote}
    Sekarang inilah asumsi baru yang akan kita gunakan untuk memperoleh
    kontradiksi. Kita memiliki dua asumsi lagi: bahwa $A \subseteq B$
    dan bahwa $B = \emptyset$. Asumsi pertama memberi kita $x \in B$:
  \end{quote}
  Karena $A \subseteq B$, maka $x \in B$.
  \begin{quote}
    Namun, karena $B = \emptyset$, setiap !!{element} dari $B$
    (misalnya $x$) juga harus merupakan !!a{element} dari~$\emptyset$.
  \end{quote}
  Karena $B = \emptyset$, maka $x \in \emptyset$. Ini merupakan
  kontradiksi, sebab menurut definisi $\emptyset$ tidak memiliki
  !!{element}.
  \begin{quote}
    Ini sudah menuntaskan bukti: dari asumsi-asumsi yang kita siapkan,
    kita telah tiba pada apa yang diperlukan (suatu kontradiksi), dan
    ini berarti semua asumsi itu tidak mungkin benar secara bersamaan.
    Karena dua asumsi pertama ($A \subseteq B$ dan $B = \emptyset$)
    tidak dipersoalkan, asumsi terakhir yang diperkenalkan (ada $x \in
    A$) pastilah salah. Namun, jika ingin benar-benar teliti, kita
    dapat menyatakannya secara eksplisit.
  \end{quote}
  Dengan demikian, asumsi kita bahwa ada $x \in A$ pastilah salah;
  jadi, melalui pembuktian dengan kontradiksi, $A$ tidak memiliki
  !!{element}.
\end{proof}

Setiap klaim positif secara sepele ekuivalen dengan suatu klaim
negatif: $p$ jika dan hanya jika $\lnot\lnot p$. Jadi, pembuktian
dengan kontradiksi juga dapat digunakan untuk menetapkan klaim positif
secara ``tidak langsung'', sebagai berikut: Untuk membuktikan~$p$,
bacalah ia sebagai klaim negatif $\lnot\lnot p$. Jika kita dapat
membuktikan suatu kontradiksi dari $\lnot p$, kita telah menetapkan
$\lnot\lnot p$ melalui pembuktian dengan kontradiksi, dan dengan
demikian menetapkan~$p$.

Dalam contoh terakhir, kita bertujuan membuktikan klaim negatif, yaitu
bahwa $A$ tidak memiliki !!{element}, sehingga asumsi yang kita buat
untuk keperluan pembuktian dengan kontradiksi (yaitu bahwa ada $x \in
A$) merupakan klaim positif. Asumsi itu memberi kita sesuatu yang
dapat diolah, yakni $x \in A$ hipotetis yang terus kita nalar sampai
kita memperoleh $x \in \emptyset$.

Ketika membuktikan klaim positif secara tidak langsung, asumsi yang
Anda buat untuk keperluan pembuktian dengan kontradiksi akan bersifat
negatif. Namun, sangat sering Anda dapat dengan mudah merumuskan
kembali klaim positif sebagai klaim negatif, dan klaim negatif sebagai
klaim positif. Bukti kita sebelumnya pada dasarnya akan tetap sama
seandainya kita membuktikan ``$A = \emptyset$'' sebagai ganti
konsekuen negatif ``$A$ tidak memiliki !!{element}.'' (Menurut definisi
$=$, ``$A = \emptyset$'' merupakan klaim umum, sebab ia terurai menjadi
``setiap !!{element} dari~$A$ merupakan !!{element} dari~$\emptyset$,
dan sebaliknya''.) Namun, pernyataan itu mudah dilihat ekuivalen dengan
klaim negatif ``tidak benar bahwa ada $x \in A$.''

Jadi, kadang-kadang lebih mudah bekerja dengan $\lnot p$ sebagai
asumsi daripada membuktikan~$p$ secara langsung. Bahkan ketika bukti
langsung sama sederhananya atau lebih sederhana (seperti dalam
contoh-contoh berikut), sebagian orang lebih suka melangkah secara
tidak langsung. Jika negasi ganda membingungkan Anda, pandanglah
pembuktian dengan kontradiksi atas suatu klaim sebagai pembuktian suatu
kontradiksi dari klaim yang \emph{berlawanan}. Jadi, pembuktian dengan
kontradiksi atas $\lnot p$ adalah pembuktian suatu kontradiksi dari
asumsi~$p$; dan pembuktian dengan kontradiksi atas~$p$ adalah
pembuktian suatu kontradiksi dari~$\lnot p$.

\begin{prop}
$A \subseteq A \cup B$.
\end{prop}

\begin{proof}
  Kita ingin menunjukkan bahwa $A \subseteq A \cup B$.
  \begin{quote}
    Sepintas lalu, ini merupakan klaim positif: setiap $x \in A$ juga
    berada dalam $A \cup B$. Negasinya ialah: terdapat suatu $x \in A$
    yang memenuhi $\notin A \cup B$. Jadi, kita dapat membuktikan
    klaim ini secara tidak langsung dengan mengasumsikan klaim yang
    dinegasikan tersebut dan menunjukkan bahwa klaim itu menuntun pada
    suatu kontradiksi.
  \end{quote}
  Andaikan tidak, yaitu $A \nsubseteq A \cup B$.
  \begin{quote}
    Kita memiliki definisi $A \subseteq A \cup B$: setiap $x \in A$
    juga memenuhi $\in A \cup B$. Untuk memahami arti $A \nsubseteq
    A \cup B$, kita harus menerapkan sedikit manipulasi logis dasar
    pada definisi yang telah diuraikan: pernyataan bahwa setiap $x \in
    A$ juga memenuhi $\in A \cup B$ adalah salah jika dan hanya jika
    ada \emph{suatu}~$x \in A$ yang memenuhi $\notin A \cup B$.
    (Di sinilah Anda harus sangat berhati-hati: banyak upaya mahasiswa
    untuk membuktikan dengan kontradiksi gagal karena mereka keliru
    menganalisis negasi klaim seperti ``semua $A$ adalah $B$''.) Dengan
    kata lain, $A \nsubseteq A \cup B$ jika dan hanya jika ada $x$
    sedemikian sehingga $x \in A$ dan $x \notin A \cup B$. Sejak itu,
    langkah selanjutnya mudah.
  \end{quote}
  Jadi, ada $x \in A$ sedemikian sehingga $x \notin A \cup B$. Menurut
  definisi $\cup$, $x \in A \cup B$ jika dan hanya jika $x \in A$ atau
  $x \in B$. Karena $x \in A$, kita memperoleh $x \in A \cup B$. Ini
  bertentangan dengan asumsi bahwa $x \notin A \cup B$.
\end{proof}

\begin{prob}
Buktikan secara \emph{tidak langsung} bahwa $A \cap B \subseteq A$.
\end{prob}

\begin{prop}
Jika $A \subseteq B$ dan $B \subseteq C$, maka $A \subseteq C$.
\end{prop}

\begin{proof}
  Andaikan $A \subseteq B$ dan $B \subseteq C$. Kita ingin menunjukkan
  $A \subseteq C$.
  \begin{quote}
    Mari kita melangkah secara tidak langsung: kita mengasumsikan
    negasi dari apa yang ingin ditetapkan.
  \end{quote}
  Andaikan tidak, yaitu $A \nsubseteq C$.
  \begin{quote}
    Seperti sebelumnya, kita menalar bahwa $A \nsubseteq C$ jika dan
    hanya jika tidak semua $x \in A$ juga memenuhi $\in C$, yaitu
    terdapat suatu $x \in A$ yang memenuhi $\notin C$. Jangan
    khawatir, dengan latihan Anda tidak perlu lagi berpikir keras untuk
    menguraikan negasi semacam ini.
  \end{quote}
  Dengan kata lain, ada~$x$ sedemikian sehingga $x \in A$ dan $x
  \notin C$.
  \begin{quote}
    Sekarang kita dapat menggunakan hal ini untuk memperoleh
    kontradiksi. Tentu saja, untuk melakukannya kita harus menggunakan
    kedua asumsi yang lain.
  \end{quote}
  Karena $A \subseteq B$, maka $x \in B$. Karena $B \subseteq C$, maka
  $x \in C$. Namun, ini bertentangan dengan $x \notin C$.
\end{proof}

\begin{prop}
Jika $A \cup B = A \cap B$, maka $A = B$.
\end{prop}

\begin{proof}
  Andaikan $A \cup B = A \cap B$. Kita ingin menunjukkan bahwa $A = B$.
  \begin{quote}
    Bagian awalnya sekarang sudah rutin:
  \end{quote}
  Untuk memperoleh kontradiksi, asumsikan bahwa $A \neq B$.
  \begin{quote}
    Asumsi kita untuk pembuktian dengan kontradiksi ialah bahwa $A
    \neq B$. Karena $A = B$ jika dan hanya jika $A \subseteq B$ dan
    $B \subseteq A$, kita memperoleh bahwa $A \neq B$ jika dan hanya
    jika $A \nsubseteq B$ \emph{atau} $B \nsubseteq A$. (Perhatikan
    betapa pentingnya berhati-hati ketika memanipulasi negasi!{}) Untuk
    membuktikan suatu kontradiksi dari disjungsi ini, kita menggunakan
    pembuktian dengan kasus dan menunjukkan bahwa dalam setiap kasus,
    suatu kontradiksi mengikuti.
  \end{quote}
  $A \neq B$ jika dan hanya jika $A \nsubseteq B$ atau $B \nsubseteq
  A$. Kita bedakan kasus.
  \begin{quote}
    Dalam kasus pertama, kita mengasumsikan $A \nsubseteq B$, yaitu
    untuk suatu $x$, $x \in A$ tetapi $\notin B$. $A \cap B$
    didefinisikan sebagai himpunan !!{element} yang dimiliki bersama
    oleh $A$ dan $B$, sehingga jika sesuatu tidak berada dalam salah
    satunya, sesuatu itu tidak berada dalam irisan. $A \cup B$ adalah
    $A$ bersama dengan $B$, sehingga segala sesuatu yang berada dalam
    salah satunya juga berada dalam gabungan. Ini memberi tahu kita
    bahwa $x \in A \cup B$, tetapi $x \notin A \cap B$, sehingga $A
    \cap B \neq A \cup B$.
  \end{quote}

  Kasus 1: $A \nsubseteq B$. Maka, untuk suatu $x$, $x \in A$, tetapi
  $x \notin B$. Karena $x \notin B$, maka $x \notin A \cap B$. Karena
  $x \in A$, maka $x \in A \cup B$. Jadi, $A \cap B \neq A \cup B$,
  bertentangan dengan asumsi bahwa $A \cap B = A \cup B$.

  Kasus 2: $B \nsubseteq A$. Maka, untuk suatu $y$, $y \in B$, tetapi
  $y \notin A$. Seperti sebelumnya, kita memperoleh $y \in A \cup B$,
  tetapi $y \notin A \cap B$, sehingga $A \cap B \neq A \cup B$,
  sekali lagi bertentangan dengan $A \cap B = A \cup B$.
\end{proof}

\end{document}
